\documentclass{ximera}

\input{../../../preamble.tex}


\definecolor{fvdnavy}{HTML}{071D43}
\definecolor{fvdcyan}{HTML}{20BFC6}
\definecolor{fvdcyanlight}{HTML}{EFFCFB}
\definecolor{fvdmagenta}{HTML}{F10D69}
\definecolor{fvdmagentalight}{HTML}{FFF1F7}
\definecolor{fvdtext}{HTML}{163044}





%=================================================
% Pedagogische omgevingen
% PDF: tcolorbox
% HTML: learning-box
%=================================================

\newenvironment{voorkennis}
{%
    \ifdefined\HCode
        \HCode{%
<section class="learning-box learning-box--prior">
<div class="learning-box__header">
<span class="learning-box__icon" aria-hidden="true"></span>
<span class="learning-box__title">Voorkennis</span>
</div>
<div class="learning-box__content">
        }%
        \begin{itemize}
    \else
       \begin{tcolorbox}[
    enhanced,
    breakable,
    colback=fvdcyanlight,
    colframe=fvdcyan,
    coltext=fvdtext,
    boxrule=0.5pt,
    leftrule=4pt,
    arc=4mm,
    outer arc=4mm,
    left=7mm,
    right=6mm,
    top=5mm,
    bottom=5mm,
    before skip=6mm,
    after skip=6mm,
    title=Voorkennis,
    coltitle=fvdcyan,
    fonttitle=\bfseries\Large,
    attach title to upper={\par\smallskip},
    boxed title style={
        empty,
        left=0mm,
        right=0mm,
        top=0mm,
        bottom=0mm
    },
    overlay={
        \draw[
            fvdcyan,
            line width=0.5pt
        ]
        ([xshift=42mm,yshift=-13mm]frame.north west)
        --
        ([xshift=-7mm,yshift=-13mm]frame.north east);
    }
]
        \begin{itemize}
    \fi
}
{%
    \end{itemize}
    \ifdefined\HCode
        \HCode{%
</div>
</section>
        }%
    \else
        \end{tcolorbox}
    \fi
}


\newenvironment{leerdoelen}
{%
    \ifdefined\HCode
        \HCode{%
<section class="learning-box learning-box--goals">
<div class="learning-box__header">
<span class="learning-box__icon" aria-hidden="true"></span>
<span class="learning-box__title">Leerdoelen</span>
</div>
<div class="learning-box__content">
        }%
        \begin{itemize}
    \else
\begin{tcolorbox}[
    enhanced,
    breakable,
    colback=fvdmagentalight,
    colframe=fvdmagenta,
    coltext=fvdtext,
    boxrule=0.5pt,
    leftrule=4pt,
    arc=4mm,
    outer arc=4mm,
    left=7mm,
    right=6mm,
    top=5mm,
    bottom=5mm,
    before skip=6mm,
    after skip=6mm,
    title=Leerdoelen,
    coltitle=fvdmagenta,
    fonttitle=\bfseries\Large,
    attach title to upper={\par\smallskip},
    boxed title style={
        empty,
        left=0mm,
        right=0mm,
        top=0mm,
        bottom=0mm
    },
    overlay={
        \draw[
            fvdmagenta,
            line width=0.5pt
        ]
        ([xshift=42mm,yshift=-13mm]frame.north west)
        --
        ([xshift=-7mm,yshift=-13mm]frame.north east);
    }
]
        \begin{itemize}
    \fi
}
{%
    \end{itemize}
    \ifdefined\HCode
        \HCode{%
</div>
</section>
        }%
    \else
        \end{tcolorbox}
    \fi
}



\addPrintStyle{../../..}

\title{Oefeningen — Implicatie en equivalentie}

\author{Ingmar Herreman}

\begin{document}

% FVD_CHAPTER_DATA_START
% Automatisch gegenereerd uit index.tex.
% Niet handmatig aanpassen.
\fvdchapterdata{1}{Logica — De taal van correct redeneren}{3}
% FVD_CHAPTER_DATA_END

\maketitle

\FVDResetExercises





\begin{oefening}[Voorwaarde en conclusie]{\oefB\ESS}

Bepaal bij elke implicatie de voorwaarde en de conclusie.

\begin{question}

Als $n$ deelbaar is door $8$, dan is $n$ even.

\begin{oplossing}

Voorwaarde:

\[
n\text{ is deelbaar door }8.
\]

Conclusie:

\[
n\text{ is even}.
\]

\end{oplossing}

\end{question}


\begin{question}

Als $x>5$, dan is $x>2$.

\begin{oplossing}

Voorwaarde:

\[
x>5.
\]

Conclusie:

\[
x>2.
\]

\end{oplossing}

\end{question}


\begin{question}

Als een driehoek gelijkzijdig is, dan heeft hij drie gelijke hoeken.

\begin{oplossing}

Voorwaarde: de driehoek is gelijkzijdig.

Conclusie: de driehoek heeft drie gelijke hoeken.

\end{oplossing}

\end{question}

\end{oefening}


% ------------------------------------------------------------

\begin{oefening}[Waarheidstabel]{\oefB\ESS}

Vul de waarheidstabel aan.

\[
\begin{array}{c|c|c}
p&q&p\Rightarrow q\\
\hline
\text{waar}&\text{waar}&\answer{\text{waar}}\\
\text{waar}&\text{onwaar}&\answer{\text{onwaar}}\\
\text{onwaar}&\text{waar}&\answer{\text{waar}}\\
\text{onwaar}&\text{onwaar}&\answer{\text{waar}}
\end{array}
\]

\begin{question}

In hoeveel rijen is de implicatie onwaar?

\begin{multipleChoice}

  \choice[correct]{In precies één rij.}

  \choice{In twee rijen.}

  \choice{In drie rijen.}

  \choice{In vier rijen.}

\end{multipleChoice}

\end{question}


\begin{question}

Beschrijf die situatie in woorden.

\begin{oplossing}

Een implicatie is onwaar wanneer de voorwaarde waar is en de conclusie
onwaar.

\end{oplossing}

\end{question}

\end{oefening}


% ------------------------------------------------------------

\begin{oefening}[Waar of onwaar?]{\oefB\ESS}

Bepaal telkens de waarheidswaarde.

\begin{question}

\[
(3<5)\Rightarrow(9<25).
\]

\begin{oplossing}

Beide deeluitspraken zijn waar.

De implicatie is dus waar.

\end{oplossing}

\end{question}


\begin{question}

\[
(7>2)\Rightarrow(7\text{ is even}).
\]

\begin{oplossing}

De voorwaarde is waar en de conclusie onwaar.

De implicatie is dus onwaar.

\end{oplossing}

\end{question}


\begin{question}

\[
(2>7)\Rightarrow(10<3).
\]

\begin{oplossing}

De voorwaarde is onwaar.

De implicatie is daarom waar.

\end{oplossing}

\end{question}


\begin{question}

\[
(2>7)\Rightarrow(10>3).
\]

\begin{oplossing}

De voorwaarde is onwaar.

De implicatie is daarom waar.

\end{oplossing}

\end{question}

\end{oefening}


% ------------------------------------------------------------

\begin{oefening}[Omgekeerde en contrapositieve]{\oefB\ESS}

Gegeven:

\[
p\Rightarrow q.
\]

\begin{question}

Welke uitspraak is de omgekeerde?

\begin{multipleChoice}

  \choice{$\neg p\Rightarrow\neg q$}

  \choice[correct]{$q\Rightarrow p$}

  \choice{$\neg q\Rightarrow\neg p$}

  \choice{$p\Leftrightarrow q$}

\end{multipleChoice}

\end{question}


\begin{question}

Welke uitspraak is de contrapositieve?

\begin{multipleChoice}

  \choice{$q\Rightarrow p$}

  \choice{$\neg p\Rightarrow\neg q$}

  \choice[correct]{$\neg q\Rightarrow\neg p$}

  \choice{$p\land q$}

\end{multipleChoice}

\end{question}


\begin{question}

Welke uitspraak heeft altijd dezelfde waarheidswaarde als
$p\Rightarrow q$?

\begin{multipleChoice}

  \choice{$q\Rightarrow p$}

  \choice[correct]{$\neg q\Rightarrow\neg p$}

  \choice{$\neg p\Rightarrow\neg q$}

\end{multipleChoice}

\end{question}

\end{oefening}




\begin{oefening}[Deelbaarheid]{\oefB\ESS}

Onderzoek voor $n\in\mathbb{Z}$ of de uitspraken waar of onwaar zijn.
Geef bij een onware uitspraak een tegenvoorbeeld.

\begin{question}

Als $n$ deelbaar is door $10$, dan is $n$ deelbaar door $5$.

\begin{oplossing}

Waar.

Als

\[
n=10k
\]

met $k\in\mathbb{Z}$, dan

\[
n=5(2k),
\]

dus is $n$ deelbaar door $5$.

\end{oplossing}

\end{question}


\begin{question}

Als $n$ deelbaar is door $5$, dan is $n$ deelbaar door $10$.

\begin{oplossing}

Onwaar.

Bijvoorbeeld

\[
n=15.
\]

Dan is $15$ deelbaar door $5$, maar niet door $10$.

\end{oplossing}

\end{question}


\begin{question}

Als $n$ deelbaar is door $12$, dan is $n$ deelbaar door $4$.

\begin{oplossing}

Waar.

Uit

\[
n=12k
\]

volgt

\[
n=4(3k).
\]

\end{oplossing}

\end{question}


\begin{question}

Als $n$ even is, dan is $n^2$ deelbaar door $4$.

\begin{oplossing}

Waar.

Als $n$ even is, bestaat er een $k\in\mathbb{Z}$ zodat

\[
n=2k.
\]

Dan

\[
n^2=(2k)^2=4k^2,
\]

dus is $n^2$ deelbaar door $4$.

\end{oplossing}

\end{question}

\end{oefening}


% ------------------------------------------------------------

\begin{oefening}[Zoek gericht een tegenvoorbeeld]{\oefU\ESS}

Onderzoek telkens de implicatie voor $x\in\mathbb{R}$.
Als ze onwaar is, geef dan een tegenvoorbeeld en leg uit waarom dit
tegenvoorbeeld de implicatie weerlegt.

\begin{question}

\[
x>5\Rightarrow x^2>25.
\]

\begin{oplossing}

Waar.

Uit $x>5$ volgt dat $x$ positief is en dus

\[
x^2>25.
\]

\end{oplossing}

\end{question}


\begin{question}

\[
x^2>25\Rightarrow x>5.
\]

\begin{oplossing}

Onwaar.

Neem bijvoorbeeld

\[
x=-6.
\]

Dan

\[
x^2=36>25,
\]

maar

\[
x>5
\]

is onwaar.

\end{oplossing}

\end{question}


\begin{question}

\[
x^2=16\Rightarrow x=4.
\]

\begin{oplossing}

Onwaar.

Neem

\[
x=-4.
\]

Dan is

\[
x^2=16
\]

waar, maar

\[
x=4
\]

onwaar.

\end{oplossing}

\end{question}


\begin{question}

\[
x=4\Rightarrow x^2=16.
\]

\begin{oplossing}

Waar.

Als $x=4$, dan volgt onmiddellijk

\[
x^2=16.
\]

\end{oplossing}

\end{question}

\end{oefening}


% ------------------------------------------------------------

\begin{oefening}[Het domein doet ertoe]{\oefU\ESS}

Beschouw de implicatie

\[
x^2=4\Rightarrow x=2.
\]

\begin{question}

Onderzoek de uitspraak voor $x\in\mathbb{R}$.

\begin{oplossing}

De uitspraak is onwaar.

Voor

\[
x=-2
\]

is de voorwaarde $x^2=4$ waar, maar de conclusie $x=2$ onwaar.

\end{oplossing}

\end{question}


\begin{question}

Onderzoek dezelfde uitspraak voor $x\in[0,+\infty[$.

\begin{oplossing}

Binnen dit domein is de uitspraak waar.

De oplossingen van

\[
x^2=4
\]

zijn $x=-2$ en $x=2$, maar $-2$ behoort niet tot het gekozen domein.

De enige toegelaten waarde waarvoor de voorwaarde waar is, is dus

\[
x=2.
\]

Dan is ook de conclusie waar.

\end{oplossing}

\end{question}


\begin{question}

Wat leert dit voorbeeld over implicaties?

\begin{oplossing}

De waarheidswaarde van een uitspraak met veranderlijken kan afhangen
van het gekozen domein.

Daarom moet het domein steeds duidelijk zijn.

\end{oplossing}

\end{question}

\end{oefening}





\begin{oefening}[Vertalen]{\oefB\ESS}

Gegeven:

\[
p\Rightarrow q.
\]

Vul telkens aan.

\begin{question}

$p$ is een \wordChoice{
\choice[correct]{voldoende}
\choice{nodige}
} voorwaarde voor $q$.

\end{question}


\begin{question}

$q$ is een \wordChoice{
\choice{voldoende}
\choice[correct]{nodige}
} voorwaarde voor $p$.

\end{question}

\end{oefening}


% ------------------------------------------------------------

\begin{oefening}[Meetkundige voorwaarden]{\oefU\ESS}

Beschouw de eigenschappen van een vierhoek:

\[
p:\quad \text{de vierhoek is een vierkant},
\]

\[
q:\quad \text{de vierhoek is een rechthoek}.
\]

\begin{question}

Welke implicatie is altijd waar?

\begin{multipleChoice}

  \choice[correct]{$p\Rightarrow q$}

  \choice{$q\Rightarrow p$}

\end{multipleChoice}

\end{question}


\begin{question}

Is vierkant zijn nodig of voldoende om rechthoek te zijn?

\begin{oplossing}

Vierkant zijn is voldoende om rechthoek te zijn.

\end{oplossing}

\end{question}


\begin{question}

Is rechthoek zijn nodig of voldoende om vierkant te zijn?

\begin{oplossing}

Rechthoek zijn is nodig om vierkant te zijn.

Het is niet voldoende, want niet elke rechthoek is een vierkant.

\end{oplossing}

\end{question}


\begin{question}

Geef een tegenvoorbeeld dat aantoont dat

\[
q\Rightarrow p
\]

onwaar is.

\begin{oplossing}

Elke rechthoek waarvan twee aangrenzende zijden verschillende lengtes
hebben is een tegenvoorbeeld.

Bijvoorbeeld een rechthoek met zijden $3$ en $5$.

\end{oplossing}

\end{question}

\end{oefening}


% ------------------------------------------------------------

\begin{oefening}[Zelf de richting bepalen]{\oefU\ESS}

Schrijf telkens de juiste implicatie en bepaal welke eigenschap nodig
en welke voldoende is.

\begin{question}

Een geheel getal is deelbaar door $20$.
Een geheel getal is deelbaar door $5$.

\begin{oplossing}

Noem

\[
p:\quad n\text{ is deelbaar door }20
\]

en

\[
q:\quad n\text{ is deelbaar door }5.
\]

Dan geldt

\[
p\Rightarrow q.
\]

Deelbaar zijn door $20$ is voldoende voor deelbaar zijn door $5$.

Deelbaar zijn door $5$ is nodig voor deelbaar zijn door $20$.

\end{oplossing}

\end{question}


\begin{question}

Een driehoek is gelijkzijdig.
Een driehoek is gelijkbenig.

\begin{oplossing}

Elke gelijkzijdige driehoek is gelijkbenig.

Dus

\[
\text{gelijkzijdig}\Rightarrow\text{gelijkbenig}.
\]

Gelijkzijdig zijn is voldoende voor gelijkbenig zijn.

Gelijkbenig zijn is nodig voor gelijkzijdig zijn.

\end{oplossing}

\end{question}

\end{oefening}





\begin{oefening}[Drie uitspraken vergelijken]{\oefU\ESS}

Beschouw voor $n\in\mathbb{Z}$:

\[
p:\quad n\text{ is deelbaar door }8,
\]

\[
q:\quad n\text{ is even}.
\]

\begin{question}

Formuleer $p\Rightarrow q$ in woorden en bepaal de waarheidswaarde.

\begin{oplossing}

Als $n$ deelbaar is door $8$, dan is $n$ even.

Deze uitspraak is waar.

\end{oplossing}

\end{question}


\begin{question}

Formuleer de omgekeerde uitspraak en bepaal de waarheidswaarde.

\begin{oplossing}

De omgekeerde is

\[
q\Rightarrow p.
\]

In woorden:

Als $n$ even is, dan is $n$ deelbaar door $8$.

Deze uitspraak is onwaar.

Bijvoorbeeld $n=6$.

\end{oplossing}

\end{question}


\begin{question}

Formuleer de contrapositieve uitspraak.

\begin{oplossing}

\[
\neg q\Rightarrow\neg p.
\]

In woorden:

Als $n$ niet even is, dan is $n$ niet deelbaar door $8$.

Voor gehele getallen:

Als $n$ oneven is, dan is $n$ niet deelbaar door $8$.

\end{oplossing}

\end{question}


\begin{question}

Welke twee van de drie uitspraken hebben noodzakelijk dezelfde
waarheidswaarde?

\begin{oplossing}

De oorspronkelijke implicatie

\[
p\Rightarrow q
\]

en haar contrapositieve

\[
\neg q\Rightarrow\neg p.
\]

\end{oplossing}

\end{question}

\end{oefening}


% ------------------------------------------------------------

\begin{oefening}[Een redenering controleren]{\oefU\ESS}

Een leerling redeneert:

\begin{center}
``Als een getal deelbaar is door $6$, dan is het deelbaar door $3$.
Het getal $15$ is deelbaar door $3$.
Dus $15$ is deelbaar door $6$.''
\end{center}

\begin{question}

Is de conclusie correct?

\begin{multipleChoice}

  \choice{Ja.}

  \choice[correct]{Nee.}

\end{multipleChoice}

\end{question}


\begin{question}

Schrijf de gegeven implicatie als $p\Rightarrow q$.

\begin{oplossing}

Neem

\[
p:\quad n\text{ is deelbaar door }6
\]

en

\[
q:\quad n\text{ is deelbaar door }3.
\]

Dan is de gegeven implicatie

\[
p\Rightarrow q.
\]

\end{oplossing}

\end{question}


\begin{question}

Welke ongeldige stap maakt de leerling?

\begin{oplossing}

De leerling redeneert van $q$ naar $p$.

Hij gebruikt dus de omgekeerde implicatie

\[
q\Rightarrow p,
\]

terwijl alleen

\[
p\Rightarrow q
\]

gegeven is.

\end{oplossing}

\end{question}

\end{oefening}




\begin{oefening}[Wanneer is er een equivalentie?]{\oefB\ESS}

Gegeven zijn twee uitspraken $p$ en $q$.

\begin{question}

Wat moet gelden om te mogen schrijven

\[
p\Leftrightarrow q?
\]

\begin{multipleChoice}

  \choice{Alleen $p\Rightarrow q$ moet waar zijn.}

  \choice{Alleen $q\Rightarrow p$ moet waar zijn.}

  \choice[correct]{Zowel $p\Rightarrow q$ als $q\Rightarrow p$ moet waar zijn.}

\end{multipleChoice}

\end{question}

\end{oefening}


% ------------------------------------------------------------

\begin{oefening}[Equivalent of niet?]{\oefU\ESS}

Onderzoek telkens of de twee eigenschappen equivalent zijn.

\begin{question}

Voor $n\in\mathbb{Z}$:

\[
p:\quad n\text{ is deelbaar door }6,
\]

\[
q:\quad n\text{ is deelbaar door }2\text{ en door }3.
\]

\begin{oplossing}

De eigenschappen zijn equivalent.

Als $n$ deelbaar is door $6$, dan is $n$ deelbaar door $2$ en door $3$.

Omgekeerd geldt voor gehele getallen: als $n$ deelbaar is door zowel
$2$ als $3$, dan is $n$ deelbaar door $6$.

Dus

\[
p\Leftrightarrow q.
\]

\end{oplossing}

\end{question}


\begin{question}

Voor $x\in\mathbb{R}$:

\[
p:\quad x=3,
\]

\[
q:\quad x^2=9.
\]

\begin{oplossing}

Niet equivalent.

Wel geldt

\[
x=3\Rightarrow x^2=9.
\]

Maar de omgekeerde uitspraak is onwaar, want $x=-3$ voldoet aan
$x^2=9$ maar niet aan $x=3$.

\end{oplossing}

\end{question}


\begin{question}

Voor $x\in[0,+\infty[$:

\[
p:\quad x=3,
\]

\[
q:\quad x^2=9.
\]

\begin{oplossing}

Binnen dit domein zijn de eigenschappen wel equivalent.

De negatieve oplossing $x=-3$ is niet toegelaten.

Dus

\[
x=3\Leftrightarrow x^2=9.
\]

\end{oplossing}

\end{question}

\end{oefening}





\begin{oefening}[Foutenanalyse]{\oefU\ESS}

Een leerling beweert:

\begin{center}
``Als $x>4$, dan is $x^2>16$.
Dus als $x^2>16$, dan is $x>4$.''
\end{center}

\begin{question}

Welke stap maakt de leerling?

\begin{oplossing}

De leerling vervangt

\[
p\Rightarrow q
\]

door de omgekeerde uitspraak

\[
q\Rightarrow p.
\]

\end{oplossing}

\end{question}


\begin{question}

Toon met een tegenvoorbeeld aan dat het besluit niet geldig is.

\begin{oplossing}

Neem bijvoorbeeld

\[
x=-5.
\]

Dan is

\[
x^2=25>16,
\]

maar

\[
x>4
\]

is onwaar.

\end{oplossing}

\end{question}


\begin{question}

Formuleer wel een uitspraak die logisch gelijkwaardig is met

\[
x>4\Rightarrow x^2>16.
\]

\begin{oplossing}

De contrapositieve is logisch gelijkwaardig:

\[
x^2\leq16\Rightarrow x\leq4.
\]

\end{oplossing}

\end{question}

\end{oefening}


% ------------------------------------------------------------

\begin{oefening}[Sensor en alarm]{\oefU\ESS}

Een laboratorium gebruikt een temperatuursensor.

De regelsoftware bevat de instructie:

\begin{center}
Als de temperatuur hoger is dan $100^\circ\mathrm{C}$,
dan wordt het alarm geactiveerd.
\end{center}

Neem

\[
p:\quad T>100^\circ\mathrm{C}
\]

en

\[
q:\quad \text{het alarm is actief}.
\]

\begin{question}

Schrijf de regel symbolisch.

\begin{oplossing}

\[
p\Rightarrow q.
\]

\end{oplossing}

\end{question}


\begin{question}

Tijdens een experiment is het alarm niet actief.
Wat mag je op basis van de gegeven regel besluiten?

\begin{oplossing}

We gebruiken de contrapositieve:

\[
\neg q\Rightarrow\neg p.
\]

Dus als het alarm niet actief is, dan is

\[
T\leq100^\circ\mathrm{C}.
\]

Dit veronderstelt uiteraard dat de opgegeven logische regel het
functioneren van het systeem correct beschrijft.

\end{oplossing}

\end{question}


\begin{question}

Tijdens een ander experiment is het alarm wel actief.
Mag je uitsluitend op basis van de gegeven regel besluiten dat

\[
T>100^\circ\mathrm{C}?
\]

\begin{oplossing}

Nee.

Dat zou de omgekeerde implicatie gebruiken:

\[
q\Rightarrow p.
\]

Die volgt niet uit de gegeven regel.

Het alarm zou eventueel ook door een andere oorzaak geactiveerd kunnen
worden.

\end{oplossing}

\end{question}

\end{oefening}


% ------------------------------------------------------------

\begin{oefening}[Een algoritme analyseren]{\oefU}

Een programma bevat de regel:

\begin{center}
Als een ingevoerd geheel getal een veelvoud is van $12$,
dan toont het programma de melding ``geschikt''.
\end{center}

Een invoer levert de melding ``geschikt'' op.

\begin{question}

Kan je uit de gegeven informatie met zekerheid besluiten dat het
ingevoerde getal een veelvoud van $12$ was?

\begin{oplossing}

Nee.

De regel zegt alleen

\[
p\Rightarrow q,
\]

waarbij

\[
p:\quad \text{het getal is een veelvoud van }12
\]

en

\[
q:\quad \text{de melding ``geschikt'' verschijnt}.
\]

Uit $q$ volgt niet automatisch $p$.

Er kunnen in het programma nog andere voorwaarden bestaan die dezelfde
melding veroorzaken.

\end{oplossing}

\end{question}


\begin{question}

Welke extra informatie zou nodig zijn om de conclusie wel te mogen
trekken?

\begin{oplossing}

We zouden ook moeten weten dat

\[
q\Rightarrow p.
\]

Dus: de melding ``geschikt'' verschijnt alleen wanneer het getal een
veelvoud van $12$ is.

Samen met de oorspronkelijke regel krijgen we dan

\[
p\Leftrightarrow q.
\]

\end{oplossing}

\end{question}

\end{oefening}





\begin{oefening}[Voorwaarden ontwerpen]{\oefG}

Een automatisch ventilatiesysteem moet aan de volgende eis voldoen:

\begin{center}
Wanneer de temperatuur hoger is dan $30^\circ\mathrm{C}$,
moet de ventilator zeker draaien.
\end{center}

Neem

\[
p:\quad T>30^\circ\mathrm{C}
\]

en

\[
q:\quad \text{de ventilator draait}.
\]

\begin{question}

Formuleer de eis symbolisch.

\begin{oplossing}

\[
p\Rightarrow q.
\]

\end{oplossing}

\end{question}


\begin{question}

Een ontwerper stelt voor om de software zo te maken dat

\[
q\Rightarrow p.
\]

Is dat dezelfde eis?

\begin{oplossing}

Nee.

De oorspronkelijke eis zegt dat een hoge temperatuur voldoende is om
de ventilator te laten draaien.

De nieuwe uitspraak zegt dat de ventilator alleen kan draaien wanneer
de temperatuur hoger is dan $30^\circ\mathrm{C}$.

Dat is een andere voorwaarde.

\end{oplossing}

\end{question}


\begin{question}

Wat betekent

\[
p\Leftrightarrow q
\]

in deze context?

\begin{oplossing}

Dan geldt zowel

\[
p\Rightarrow q
\]

als

\[
q\Rightarrow p.
\]

De ventilator draait dus precies wanneer de temperatuur hoger is dan
$30^\circ\mathrm{C}$.

Een temperatuur boven $30^\circ\mathrm{C}$ is dan zowel nodig als
voldoende voor het draaien van de ventilator.

\end{oplossing}

\end{question}

\end{oefening}


% ------------------------------------------------------------

\begin{oefening}[Een bewering onderzoeken]{\oefG}

Een leerling formuleert voor $x\in\mathbb{R}$:

\[
x^2>1\Rightarrow x>1.
\]

Een tweede leerling zegt:

\begin{center}
``De uitspraak is bijna juist. We moeten alleen $>$ vervangen door
$\geq$.''
\end{center}

\begin{question}

Onderzoek de oorspronkelijke uitspraak.

\begin{oplossing}

Ze is onwaar.

Bijvoorbeeld

\[
x=-2
\]

geeft

\[
x^2=4>1,
\]

maar

\[
x>1
\]

is onwaar.

\end{oplossing}

\end{question}


\begin{question}

Onderzoek het voorstel

\[
x^2\geq1\Rightarrow x\geq1.
\]

\begin{oplossing}

Ook deze uitspraak is onwaar.

Opnieuw is bijvoorbeeld

\[
x=-2
\]

een tegenvoorbeeld.

\end{oplossing}

\end{question}


\begin{question}

Formuleer zelf een correcte conclusie $q$ zodat

\[
x^2>1\Rightarrow q
\]

voor alle $x\in\mathbb{R}$ waar is en de conclusie zo precies mogelijk
de mogelijke waarden van $x$ beschrijft.

\begin{oplossing}

Uit

\[
x^2>1
\]

volgt

\[
|x|>1.
\]

Dit betekent

\[
x<-1\quad\text{of}\quad x>1.
\]

Een correcte implicatie is dus

\[
x^2>1
\Rightarrow
(x<-1\lor x>1).
\]

Sterker nog: beide voorwaarden zijn equivalent:

\[
x^2>1
\Leftrightarrow
(x<-1\lor x>1).
\]

\end{oplossing}

\end{question}

\end{oefening}


% ------------------------------------------------------------

\begin{oefening}[Zelf een tegenvoorbeeld construeren]{\oefG}

Bedenk zelf voor elk van de volgende situaties een voorbeeld van twee
wiskundige eigenschappen $p$ en $q$.

\begin{question}

$p\Rightarrow q$ is waar, maar $q\Rightarrow p$ is onwaar.

\begin{oplossing}

Er zijn veel mogelijkheden.

Bijvoorbeeld voor $n\in\mathbb{Z}$:

\[
p:\quad n\text{ is deelbaar door }4,
\]

\[
q:\quad n\text{ is even}.
\]

Dan is

\[
p\Rightarrow q
\]

waar, maar

\[
q\Rightarrow p
\]

onwaar, bijvoorbeeld voor $n=6$.

\end{oplossing}

\end{question}


\begin{question}

$p\Rightarrow q$ en $q\Rightarrow p$ zijn beide waar.

\begin{oplossing}

Bijvoorbeeld voor $n\in\mathbb{Z}$:

\[
p:\quad n\text{ is even},
\]

\[
q:\quad n^2\text{ is even}.
\]

Dan geldt

\[
p\Leftrightarrow q.
\]

Andere correcte voorbeelden zijn mogelijk.

\end{oplossing}

\end{question}


\begin{question}

Leg uit waarom de tweede situatie precies betekent dat $p$ en $q$
equivalent zijn.

\begin{oplossing}

Een equivalentie

\[
p\Leftrightarrow q
\]

betekent dat beide richtingen waar zijn:

\[
p\Rightarrow q
\]

en

\[
q\Rightarrow p.
\]

\end{oplossing}

\end{question}

\end{oefening}



\end{document}


